\begin{abstract}
We develop a cybernetic model of economic planning that
integrates dynamic input-output analysis, endogenous growth, and
shadow-price-guided coordination into a unified optimization framework.
Growth is inferred from revealed consumption demand rather than imposed
exogenously, while capital accumulation is embedded directly into the
technical production structure.
A constrained social utility maximization problem generates shadow
prices that guide decentralized price adjustment and adaptive preference
evolution.
The resulting system operates through two nested feedback loops: a fast
price-and-inventory adjustment loop and a slow structural growth loop.
We benchmark the system against a Heterogeneous Agent New Keynesian (HANK) 
model over a 20-quarter horizon using 2022 Spanish input-output data. 
Monte Carlo results demonstrate that while the HANK model produces 
higher raw GDP growth (~15.5\% vs 12.3\%), the cybernetic planner 
achieves an 8.3\% welfare advantage by eliminating monetary 
distortions and uncoordinated over-investment, achieving superior 
allocative efficiency.
This framework offers a computationally tractable alternative to market
coordination and classical material balance planning, partially addressing key
critiques raised by \citet{hayek1945use} while incorporating modern
computational methods unavailable to early socialist planning theorists
\citep{lange1936economic,taylor1929guidance}.

\medskip
\noindent\textbf{JEL Codes:}  C67, D58, P21, P27

\noindent\textbf{Keywords:} Economic Planning; Input-Output Analysis;
Shadow Prices; Cybernetics; Computational Economics
\end{abstract}